% !TEX program = xelatex
% AI-effects 프로젝트 — Anthropic Economic Index(AEI)가 HuggingFace에 공개하는
% 데이터셋 저장소(Anthropic/EconomicIndex) 자체를 직접 조회해, (1) 과업의 Claude
% 사용/미사용 노동시간 데이터, (2) 과업·직업별 사용량 데이터의 필드 정의·측정단위를
% 정리하고, 선행 노출도(exposure)·적용률(adoption) 개념과의 접점을 논의한 작업 문헌(report/).
% 컴파일: xelatex -> biber -> xelatex -> xelatex (kotex, biblatex/biber 필요)
\documentclass[11pt]{article}

\usepackage{fontspec}
\usepackage{kotex}
\usepackage[a4paper,margin=2.2cm]{geometry}
\usepackage{longtable}
\usepackage{booktabs}
\usepackage{array}
\usepackage{amsmath}
\usepackage{xcolor}
\usepackage{hyperref}
\usepackage{enumitem}
\usepackage{titlesec}
\usepackage[backend=biber,style=authoryear,maxcitenames=2,uniquelist=false,sorting=nyt]{biblatex}
\addbibresource{../../reference/references.bib}

\hypersetup{colorlinks=true,linkcolor=blue!50!black,citecolor=blue!50!black,urlcolor=blue!50!black}
\newcolumntype{L}[1]{>{\raggedright\arraybackslash}p{#1}}
\setlength{\LTleft}{0pt}
\setlength{\LTright}{0pt}
\renewcommand{\arraystretch}{1.15}
\setlength{\tabcolsep}{4pt}

\title{Anthropic Economic Index 데이터 해부:\\[1mm]
과업 노동시간과 과업·직업별 사용량\\[2mm]
\large --- Claude 사용/미사용 시간 데이터와 과업·직업별 사용량 데이터의 정의, 측정 단위, 그리고 선행 노출도·적용률 개념과의 접점 ---}
\author{AI-effects 프로젝트}
\date{\today}

\begin{document}
\maketitle

\begin{abstract}
\noindent 본 문헌은 \texttt{07\_AEI실증연구\_노출도사용량\_계보.tex}가 다룬 개별 실증연구들이 아니라, 그 연구들이 공통으로 사용하는 \textbf{Anthropic Economic Index(AEI) 데이터셋 자체}의 구조를 직접 해부한다. Anthropic이 HuggingFace에 공개하는 \texttt{Anthropic/EconomicIndex} 저장소(6개 릴리스 폴더+\texttt{labor\_market\_impacts} 폴더, 각 릴리스의 \texttt{data\_documentation.md}/\texttt{README.md})를 2026년 9월 시점 스냅샷으로 직접 조회해, (1) \textbf{과업의 Claude 사용/미사용 시 노동시간 데이터}(\texttt{human\_only\_time}·\texttt{human\_with\_ai\_time} 등 "경제원시지표" 계열 필드)와 (2) \textbf{과업·직업별 사용량 데이터}(\texttt{onet\_task}·\texttt{soc\_occupation}·\texttt{request} 등 facet의 건수·비중)의 정의, 측정 단위, 코딩 방식, 도입 시점을 필드 단위로 정리한다. 나아가 이 두 데이터가 선행 AI 노출도(exposure) 연구---Felten et al.(2021/2023), Webb(2020), Eloundou et al.(2023)---및 AI 적용률(adoption) 연구---기업·근로자 자기보고 서베이, 전병유·신영민(2025)의 AI Firm Exposure(AIFE)---의 개념과 각각 어떻게 연결·대비되는지를 논의한다. 핵심 결론은 다음과 같다: (2) 사용량 데이터는 "예측된 노출"을 "관측된 사용"으로 대체하는 데이터이고, (1) 노동시간 데이터는 노출도 지표가 이분법적·서수형으로만 표현하던 "잠재력"을 시간절감률이라는 \textbf{연속적 실현 크기}로 정량화하는 데이터다. 두 데이터는 \texttt{labor\_market\_impacts/job\_exposure.csv}·\texttt{task\_penetration.csv} 두 파일에서 이미 하나의 지수("관측노출")로 결합되어 있으며, 이는 본 프로젝트가 \texttt{data/proc/merged/}에서 재현하고자 하는 결합 구조의 실제 선례다.
\end{abstract}

\tableofcontents

% ============================================================
\section{서론}
% ============================================================

\subsection{목적과 기존 report와의 관계}

\texttt{07\_AEI실증연구\_노출도사용량\_계보.tex}는 AEI 데이터를 활용한 실증연구 9편이 각각 "노출도"와 "사용량"을 어떻게 \emph{논의}하는지를 다뤘다. 그러나 그 논의는 모두 \textbf{AEI가 발행하는 원자료 자체가 정확히 어떤 필드로 구성되어 있는지}를 전제로 한다. 본 문헌은 그 전제를 직접 검증하기 위해, Anthropic이 실제로 공개한 데이터셋 저장소의 문서화 파일을 조회해 필드 단위로 재구성한 자료 해설서다. 목적은 두 가지다.

\begin{enumerate}[label=(\arabic*)]
  \item 향후 본 프로젝트가 AEI 원자료(HuggingFace CSV)를 직접 다운로드해 \texttt{data/raw/usage/}에 적재할 경우, 어떤 열(column)을 어떤 정의$\cdot$단위로 읽어야 하는지 사전에 문서화한다.
  \item (1) 노동시간 데이터와 (2) 사용량 데이터가 선행 노출도$\cdot$적용률 개념과 정확히 어느 지점에서 대응하는지를, 필드 수준의 근거로 논증한다.
\end{enumerate}

\subsection{자료 출처와 조회 방법}

Anthropic Economic Index의 데이터는 \url{https://huggingface.co/datasets/Anthropic/EconomicIndex}에 공개되어 있다. 이 저장소는 2026년 9월 시점 다음과 같은 구조를 갖는다(저장소 최상위 \texttt{README.md}와 각 릴리스 폴더의 \texttt{data\_documentation.md}/\texttt{README.md}를 직접 조회해 확인).

\begin{longtable}{L{2.6cm}L{3.0cm}L{5.0cm}L{5.5cm}}
\caption{\texttt{Anthropic/EconomicIndex} 저장소 구조 개관}\label{tab:repo}\\
\toprule
\textbf{폴더} & \textbf{대응 보고서(발행일)} & \textbf{스키마 형식} & \textbf{비고}\\
\midrule
\endfirsthead
\multicolumn{4}{l}{\small (표 \thetable\ 계속)}\\
\toprule
\textbf{폴더} & \textbf{대응 보고서(발행일)} & \textbf{스키마 형식} & \textbf{비고}\\
\midrule
\endhead
\bottomrule
\endfoot

\texttt{release\_2025\_02\_10} &
\citet{handa-2025-which-economic-tasks-are}, R1 &
개별 flat CSV(\texttt{onet\_task\_mappings.csv}, \texttt{automation\_vs\_augmentation.csv} 등) &
과업$\cdot$협업패턴별로 파일이 분리된 초기 형태\\

\texttt{release\_2025\_03\_27} &
(R2, 국가분류 없는 세계 전체 웨이브) &
flat CSV+\texttt{cluster\_level\_data/} 하위폴더 & \\

\texttt{release\_2025\_09\_15} &
\citet{appel-2025-uneven-geographic-and-enterprise-ai}, V3 &
장방형(long-format) 통합 패널 최초 도입: \texttt{geo\_id}/\texttt{geography}/\texttt{facet}/\texttt{level}/\texttt{variable}/\texttt{cluster\_name}/\texttt{value} &
지리적 분해(AUI) 최초 도입 웨이브이나, AUI 자체는 저장된 열이 아니라 사용점유율$\div$인구점유율로 논문이 직접 계산하는 파생지표\\

\texttt{release\_2026\_01\_15} &
\citet{appel-2026-economic-primitives}, V4 &
위와 동일 스키마+\texttt{human\_only\_time}/\texttt{human\_with\_ai\_time} 등 "경제원시지표" facet 신규 추가 &
부록 \texttt{aei\_v4\_appendix.pdf}에 9개 분류기 프롬프트 전문 수록\\

\texttt{release\_2026\_03\_24} &
\citet{massenkoff-2026-learning-curves}, V5 &
동일 스키마 유지 & \\

\texttt{release\_2026\_06\_26} &
\citet{massenkoff-2026-cadences}, V6 &
열 이름 재정의: \texttt{facet}$\to$\texttt{category\_name}, \texttt{variable}$\to$\texttt{metric\_id}, \texttt{cluster\_name}$\to$\texttt{node\_name}, \texttt{level}$\to$\texttt{hierarchy\_level} &
산출물(artifact) 분류 24종 신규 추가\\

\texttt{labor\_market\_impacts} &
\citet{massenkoff-2026-labor-market-impacts-of-ai} &
독립 CSV 2종: \texttt{job\_exposure.csv}, \texttt{task\_penetration.csv} &
(1)과 (2)를 결합한 "관측노출" 산출 파생 데이터\\
\end{longtable}

이 표가 보여주듯, AEI의 데이터 스키마는 고정되어 있지 않고 릴리스마다 진화해 왔다---R1의 목적별 개별 파일 구조에서 V3부터 장방형 통합 패널로, V6에서 다시 열 이름이 재정의되는 식이다. 이는 AEI 데이터를 \texttt{data/proc/}에 흡수할 때 \textbf{릴리스별로 파싱 로직을 분리}해야 함을 시사한다.

% ============================================================
\section{(1) 과업 노동시간 데이터: 정의와 측정 단위}
\label{sec:time}
% ============================================================

\subsection{도입 배경: Tamkin and McCrory(2025)에서 Appel et al.(2026)로}

노동시간 데이터는 원래 AEI 공식 시리즈의 일부가 아니라, \citet{tamkin-2025-estimating-ai-productivity-gains-from}가 별도 응용연구로 개발한 방법론이었다. 이들은 Claude.ai 대화 10만 건 각각에 대해 Claude 자신에게 (i) "AI 지원 없이 전문가가 이 과업을 완료하는 데 걸리는 시간"과 (ii) "AI 지원으로 실제 걸린 시간"을 프롬프트로 직접 추정시켰다---전자는 시간(hour) 단위로, 후자는 분(minute) 단위로 응답하도록 설계되었다(부록 프롬프트 "Human time estimation prompt"/"Interaction time estimation prompt"). 이 방법론은 자기일관성 검정(로그상관 $r=0.89$--$0.93$)과 외부 벤치마크(JIRA 티켓, Spearman $\rho=0.44$--$0.50$)로 검증되었다.

이 방법론은 이후 \citet{appel-2026-economic-primitives}(V4 "경제원시지표")에서 AEI 공식 데이터셋의 정식 facet으로 흡수되었다---\texttt{release\_2026\_01\_15}부터 \texttt{human\_only\_time}$\cdot$\texttt{human\_with\_ai\_time} facet이 최초로 등장하며, 직전 릴리스(\texttt{release\_2025\_09\_15}, V3)의 \texttt{data\_documentation.md}에는 이 facet이 \textbf{존재하지 않음}을 직접 확인했다. 즉 "노동시간 데이터"는 AEI 시리즈 전체가 아니라 \textbf{V4(2026년 1월) 이후}에만 공식 데이터로 존재한다.

\subsection{필드 정의}

\begin{longtable}{L{2.6cm}L{5.4cm}L{2.2cm}L{5.4cm}}
\caption{과업 노동시간 관련 facet 정의(\texttt{release\_2026\_01\_15} \texttt{data\_documentation.md} 기준)}\label{tab:time}\\
\toprule
\textbf{facet 이름} & \textbf{정의} & \textbf{도입 시점} & \textbf{비고}\\
\midrule
\endfirsthead
\multicolumn{4}{l}{\small (표 \thetable\ 계속)}\\
\toprule
\textbf{facet 이름} & \textbf{정의} & \textbf{도입 시점} & \textbf{비고}\\
\midrule
\endhead
\bottomrule
\endfoot

\texttt{human\_only\_time} &
사람이 AI 지원 없이 해당 과업을 완료하는 데 필요한 추정 시간 &
V4(2026.1) &
방법론상 시간(hour) 단위로 프롬프트 설계되었으나, 공개 데이터셋 문서에는 단위가 명시되어 있지 않음(\S\ref{sec:time-caveat} 참조)\\

\texttt{human\_with\_ai\_time} &
사람이 AI 지원을 받아 해당 과업을 완료하는 데 걸린 추정 시간 &
V4(2026.1) &
위와 동일 사유로 방법론상 분(minute) 단위이나, 문서상 단위 미명시\\

\texttt{human\_only\_ability} &
사람이 AI 없이 해당 과업을 단독으로 완료할 수 있는지 여부 &
V4(2026.1) &
범주형(완료가능/완료불가 등); \citet{massenkoff-2026-learning-curves}가 인용한 "인간단독완료불가 비율"(15.35\%$\to$14.30\%, 2025.11$\to$2026.2)의 원 출처\\

\texttt{task\_success} &
해당 과업(대화)이 성공적으로 완료되었는지 여부 &
V4(2026.1) &
\citet{appel-2026-economic-primitives}의 "실효 AI 커버리지" 산출에 사용되는 성공률의 원 필드\\

\texttt{ai\_autonomy} &
과업 수행 중 AI가 스스로 결정을 내리는 정도 &
V4(2026.1) &
1(전무)$\sim$5(매우 높음) 서열 척도(원 보고서 서술 기준)\\

\texttt{human\_education\_years} &
프롬프트를 이해하는 데 필요한 사람의 교육연수 추정치 &
V4(2026.1) &
\citet{appel-2026-economic-primitives}의 순탈숙련(net deskilling) 분석에 사용\\

\texttt{ai\_education\_years} &
Claude의 응답을 이해하는 데 필요한 교육연수 추정치(응답의 "숙련도") &
V4(2026.1) &
위와 동일 분석에 사용, 인간-AI 교육연수 상관 $r=0.925$(국가수준)\\

\texttt{use\_case} &
대화의 용도 구분(업무/코스워크/개인) &
V4(2026.1) &
\citet{massenkoff-2026-cadences}(V6)의 산출물-용도 교차분석에도 재사용\\
\end{longtable}

이 8개 facet은 모두 \textbf{개별 대화(conversation) 단위로 Claude 자신이 직접 분류$\cdot$추정한 값}을 O*NET 과업$\cdot$직업 또는 지리 단위로 집계한 것이다(\citealp{appel-2026-economic-primitives}, "9개 신규 분류기"). 검증은 인간 연구자와의 비교(소표본) 및 방향적 정확성(directional accuracy) 확인에 그치며, 개별 값의 정밀도는 보장되지 않음이 원 보고서에 명시되어 있다.

\subsection{파일 구조와 집계 통계량}

\texttt{release\_2026\_01\_15}의 장방형 스키마에서, 위 facet들은 다음과 같은 \texttt{variable}(집계통계량) 값을 갖는다: \texttt{mean}, \texttt{median}, \texttt{stdev}, \texttt{mean\_ci\_lower}, \texttt{mean\_ci\_upper}, \texttt{median\_ci\_lower}, \texttt{median\_ci\_upper}, \texttt{count}, \texttt{histogram\_count}, \texttt{histogram\_pct}. 즉 하나의 과업(\texttt{cluster\_name})에 대해 "평균 소요시간이 얼마인가"뿐 아니라 "그 분포의 신뢰구간$\cdot$히스토그램"까지 함께 제공되어, 단일 평균값이 아니라 분포 전체를 이용한 분석(예: \citealp{appel-2026-economic-primitives}의 speedup 분포 비교)이 가능하도록 설계되어 있다.

시간절감률(time savings)은 저장된 필드가 아니라, \texttt{human\_only\_time}과 \texttt{human\_with\_ai\_time}의 평균값으로부터 연구자가 직접 계산하는 \textbf{파생 지표}다:
\[
\text{time savings} = 1 - \frac{\text{human\_with\_ai\_time}}{\text{human\_only\_time}}.
\]
\citet{tamkin-2025-estimating-ai-productivity-gains-from}는 이를 Hulten 정리로 집계해 거시 노동생산성 증가율(+1.8\%p)을 산출했다.

\subsection{유의사항: 측정 단위의 문서화 공백}
\label{sec:time-caveat}

\texttt{data\_documentation.md}를 직접 조회한 결과, \textbf{\texttt{human\_only\_time}$\cdot$\texttt{human\_with\_ai\_time}의 측정 단위(시간 vs 분)는 데이터셋 문서 어디에도 명시되어 있지 않다.} \citet{tamkin-2025-estimating-ai-productivity-gains-from}의 방법론 절(원 논문 p.5)은 "AI 지원 없이... 시간단위로", "AI 지원으로... 분단위로" 프롬프트를 설계했다고 서술하지만, 이 서술이 V4 공식 데이터셋에 그대로 반영되었는지---즉 두 facet이 서로 다른 단위로 저장되어 있는지, 아니면 공식 데이터셋화 과정에서 단일 단위(예: 분)로 통일되었는지---는 문서만으로 확인할 수 없다. \textbf{본 프로젝트가 이 데이터를 \texttt{data/proc/}에 적재할 경우, 반드시 원자료 CSV의 실제 표본값(예: 소프트웨어 개발 관련 과업의 \texttt{human\_only\_time} 평균값이 1$\sim$10 범위인지 60$\sim$600 범위인지)을 직접 확인해 단위를 역으로 검증하는 절차를 do 파일 초반에 넣어야 한다}---이는 재현성 원칙(\texttt{CLAUDE.md})상 반드시 \texttt{01\_import\_*.do}의 주석으로 남겨야 할 사안이다.

% ============================================================
\section{(2) 과업·직업별 사용량 데이터: 정의와 측정 단위}
\label{sec:usage}
% ============================================================

\subsection{원형: R1(Handa et al. 2025)의 flat CSV}

R1(\texttt{release\_2025\_02\_10})의 사용량 데이터는 목적별로 분리된 개별 CSV로 존재한다.

\begin{itemize}[leftmargin=*]
  \item \texttt{onet\_task\_mappings.csv}: 열 \texttt{task\_name}(과업 기술문), \texttt{pct}(해당 과업을 포함하는 대화의 비율)로 구성---O*NET 약 2만 개 과업문 중 Clio가 각 대화를 매핑한 결과를 과업 단위로 집계한 것이다.
  \item \texttt{automation\_vs\_augmentation.csv}: 열 \texttt{interaction\_type}(directive/feedback loop/task iteration/learning/validation 5유형), \texttt{pct}(해당 패턴을 보이는 대화의 비율)로 구성---협업모드 분류 결과다.
  \item 이 외 \texttt{SOC\_Structure.csv}(직업 대분류 체계), \texttt{bls\_employment\_may\_2023.csv}(BLS 고용통계), \texttt{onet\_task\_statements.csv}(O*NET 과업 원문), \texttt{wage\_data.csv}(임금자료)가 매칭$\cdot$가중치 계산용 보조 테이블로 함께 제공된다.
\end{itemize}

\subsection{V3 이후: 장방형 통합 패널}

\texttt{release\_2025\_09\_15}(V3)부터는 사용량 데이터가 \texttt{human\_only\_time} 계열과 동일한 장방형 스키마(\texttt{geo\_id}/\texttt{geography}/\texttt{facet}/\texttt{level}/\texttt{variable}/\texttt{cluster\_name}/\texttt{value})로 통합된다. 사용량과 직접 관련된 facet은 다음과 같다.

\begin{longtable}{L{2.4cm}L{5.4cm}L{3.6cm}L{4.2cm}}
\caption{과업·직업별 사용량 관련 facet 정의}\label{tab:usagefacet}\\
\toprule
\textbf{facet 이름} & \textbf{정의} & \textbf{level(0--2) 의미} & \textbf{도입 시점}\\
\midrule
\endfirsthead
\multicolumn{4}{l}{\small (표 \thetable\ 계속)}\\
\toprule
\textbf{facet 이름} & \textbf{정의} & \textbf{level(0--2) 의미} & \textbf{도입 시점}\\
\midrule
\endhead
\bottomrule
\endfoot

\texttt{onet\_task} &
O*NET 과업 단위 사용 비중/건수 &
0=개별 과업, 상위레벨=과업군 &
R1(2025.2)부터, V3부터 장방형 스키마로 통합\\

\texttt{soc\_occupation} &
SOC(Standard Occupational Classification) 직업 단위 사용 비중/건수 &
0=세부직업(6자리), 상위레벨=대분류(2자리) &
V3(2025.9)부터 명시적 facet화\\

\texttt{request} &
요청 내용 기반 군집(직업 분류와 독립적인 "요청군집") &
0=최하위 세부과업, 상위=대분류 &
V3(2025.9)부터, \citet{fan-2026-what-countries-use-ai-and}이 "두 개의 렌즈" 중 하나로 채택\\

\texttt{collaboration} &
협업모드(자동화 대 증강) 5유형 분류 &
단일 레벨 &
R1(2025.2)부터(구 \texttt{automation\_vs\_augmentation.csv})\\

\texttt{multitasking} &
한 대화 내 여러 과업 동시 수행 여부 &
단일 레벨 &
V3 이후\\
\end{longtable}

각 facet의 \texttt{variable}(집계통계량)은 \texttt{count}(해당 셀에 매핑된 대화 건수)와 \texttt{pct}(비율)가 기본이며, \texttt{pct}의 정의는 데이터 문서에 "부모 지리 단위 대비 전체 사용량의 비율(국가는 전세계 대비, 미국 주는 소속 국가 대비)"로 명시되어 있다---즉 \textbf{분모가 항상 상위 지리 단위의 총사용량}이라는 점이 R1의 flat CSV(분모가 항상 전체 표본)와 구별되는 V3 이후의 특징이다.

\subsection{"두 개의 렌즈": 직업 렌즈 대 요청 렌즈}

\citet{fan-2026-what-countries-use-ai-and}이 명시적으로 정리했듯, AEI의 사용량 데이터는 대화를 두 가지 독립적인 방식으로 분류한다.

\begin{enumerate}[label=(\arabic*)]
  \item \textbf{직업 렌즈}(\texttt{onet\_task}$\to$\texttt{soc\_occupation}): 미국 노동시장 데이터(O*NET)로 훈련된 분류기가 대화를 직업 과업에 매핑. 미국 중심 분류체계이므로 비미국 국가에서는 미분류(not-classified) 비율이 훨씬 높다(\citealp{fan-2026-what-countries-use-ai-and} 보고: 국가수준 미분류 비율 중앙값이 직업 렌즈 66\% vs 요청 렌즈 7\%).
  \item \textbf{요청 렌즈}(\texttt{request}): 직업 분류와 무관하게 대화의 내용 자체로 군집화(3단계 위계, 621개 최하위 세부과업$\to$104개 level-1$\to$26개 level-2 대분류). 국가간 비교에는 이 렌즈가 더 안정적이라는 것이 \citet{fan-2026-what-countries-use-ai-and}의 판단이다.
\end{enumerate}

이는 "과업$\cdot$직업별 사용량"이라는 단일한 개념이 실제로는 \textbf{미국 노동시장 프레임에 고정된 지표}(직업 렌즈)와 \textbf{프레임에 중립적인 지표}(요청 렌즈)로 나뉜다는 뜻이며, 본 프로젝트가 한국 노동시장과 AEI 데이터를 결합할 때 직업 렌즈를 그대로 쓸 경우 미분류율이 높을 수 있음을 시사한다(\S\ref{sec:conclusion}에서 재론).

\subsection{지리적 분해와 프라이버시 임계치}

\texttt{geo\_id}는 R1$\sim$V2 시기에는 존재하지 않다가, V3(\texttt{release\_2025\_09\_15})부터 도입되었다. 형식은 국가의 경우 ISO-3166-1(V3는 2자리, V4부터는 3자리로 변경), 미국 주의 경우 ISO 3166-2, 그 외 \texttt{"GLOBAL"}이다. 셀 포함 여부에는 두 층위의 프라이버시 임계치가 적용된다.

\begin{itemize}[leftmargin=*]
  \item \textbf{국가/지역 단위 포함 임계치}: 국가당 최소 200건 대화, 미국 주당 최소 100건 대화("Minimum Observations", V3$\cdot$V4 문서에 동일하게 명시)---이 기준 미만인 지역은 데이터셋에서 아예 제외된다("enrichment 단계에서 적용, raw 전처리 단계 아님"이라는 단서도 명시).
  \item \textbf{셀(cell) 단위 억제 임계치}: \citet{appel-2025-uneven-geographic-and-enterprise-ai}의 원 보고서에서 확인한 대로, 국가-과업 셀은 최소 15건 대화$\cdot$5개 고유 계정, 하위(bottom-up) 요청군집은 최소 500건 대화$\cdot$250개 고유 계정을 충족해야 한다.
\end{itemize}

두 임계치는 서로 다른 단계에 적용되는 별개의 기준이다(전자는 "이 국가를 데이터셋에 포함할지", 후자는 "이 국가 내에서 이 특정 과업 셀을 노출할지").

\subsection{파생 지표: AUI}

\citet{appel-2025-uneven-geographic-and-enterprise-ai}이 정의한 AUI(Anthropic AI Usage Index, 지역 사용점유율$\div$지역 근로연령인구점유율)는 \textbf{저장된 열이 아니다.} 데이터셋 문서를 직접 확인한 결과 AUI라는 이름의 필드는 어느 릴리스에도 존재하지 않으며, 이는 \texttt{soc\_occupation}/\texttt{request} facet의 \texttt{pct} 값과 외부 인구통계(UN/World Bank 근로연령인구 자료)를 연구자가 직접 결합해 사후에 계산하는 지수다. 마찬가지로 \citet{fan-2026-what-countries-use-ai-and}의 확산폭(로그HHI)$\cdot$강도(PPML 회귀 대상), \citet{fan-2026-aggregate-gains-from-ai}의 LCE$\cdot$ACI도 모두 원자료의 \texttt{count}/\texttt{pct} 값에 외부 임금$\cdot$인구 자료를 결합해 연구자가 사후 계산하는 2차 지표이지, AEI가 직접 제공하는 필드가 아니다.

% ============================================================
\section{(1)과 (2)의 결합: \texttt{labor\_market\_impacts} 폴더}
\label{sec:combine}
% ============================================================

AEI 저장소에서 (1) 노동시간 데이터와 (2) 사용량 데이터가 실제로 결합되어 있는 유일한 지점은 \texttt{labor\_market\_impacts} 폴더다. 이 폴더는 릴리스 날짜별 하위구조 없이 독립적으로 존재하며, 단 2개의 CSV만 포함한다.

\begin{longtable}{L{3.0cm}L{5.5cm}L{7.3cm}}
\caption{\texttt{labor\_market\_impacts} 폴더 구성}\label{tab:lmi}\\
\toprule
\textbf{파일} & \textbf{내용(추정 근거)} & \textbf{\citet{massenkoff-2026-labor-market-impacts-of-ai} 원 논문과의 대응}\\
\midrule
\endfirsthead
\multicolumn{3}{l}{\small (표 \thetable\ 계속)}\\
\toprule
\textbf{파일} & \textbf{내용(추정 근거)} & \textbf{원 논문과의 대응}\\
\midrule
\endhead
\bottomrule
\endfoot

\texttt{task\_penetration.csv} &
과업 단위 사용량(usage) 침투도 &
원 논문의 $WorkUsage_t$(과업 $t$의 가중 업무사용량)에 대응---\S\ref{sec:time}$\cdot$\S\ref{sec:usage}의 사용량 데이터를 과업 단위로 정리한 것\\

\texttt{job\_exposure.csv} &
직업 단위 "관측노출(observed exposure)" 지수 $R_o$ &
Eloundou $\beta$(이론노출)를 $WorkUsage_t \geq 100$ 게이트로 걸러 자동화비중 가중치 $\alpha_t$로 집계한 최종 지수---\S\ref{sec:time}의 성공/자동화 관련 필드와 \S\ref{sec:usage}의 사용량 필드를 모두 사용해 산출된 \textbf{결합형 지표}의 공개 배포판\\
\end{longtable}

이 구조는 \texttt{07\_AEI실증연구\_노출도사용량\_계보.tex}가 논증한 "9편 중 유일한 계산적 결합 사례"가 실제로 데이터셋 저장소에서도 \textbf{별도 폴더로 격리되어 배포}되고 있음을 보여준다---즉 Anthropic 스스로도 "순수 사용량 데이터"(릴리스별 폴더)와 "노출$\times$사용 결합 데이터"(\texttt{labor\_market\_impacts})를 구조적으로 분리해 취급하고 있다.

% ============================================================
\section{선행연구의 AI 노출도·AI 적용률 개념과의 연결}
\label{sec:link}
% ============================================================

\subsection{(2) 사용량 데이터 $\leftrightarrow$ AI 노출도(exposure) 개념: 예측에서 관측으로}

\texttt{02\_AI노출도\_지표\_비교.tex}가 정리한 \citet{felten-2021-occupational-industry-and-geographic-exposure}의 AIOE, \citet{webb-2020-impact-artificial-intelligence-labor}의 특허기반 지수, \citet{eloundou-2023-gpts-are-gpts-an-early}의 GPT 노출 루브릭은 모두 \textbf{직업$\cdot$과업 단위로 "AI가 그 일을 할 수 있는가"를 사전에 추정}한 스칼라 값이다---크라우드소싱된 인간 평가자, 특허텍스트의 어휘적 유사도, 또는 LLM 평가자가 그 값을 산출한다는 점에서 방법은 다르지만, 공통적으로 \textbf{"관측"이 아니라 "판정"}이다.

AEI의 \texttt{onet\_task}/\texttt{soc\_occupation}/\texttt{request} facet은 이와 근본적으로 다른 종류의 숫자다. \texttt{pct}$\cdot$\texttt{count} 값은 판정이 아니라 \textbf{이미 일어난 대화의 집계}이며, 따라서 이론적으로는 "노출 가능"하지만 실제로는 전혀 사용되지 않는 과업(\citealp{massenkoff-2026-labor-market-impacts-of-ai}가 발견한, Computer \& Math 직업군에서 이론노출 94\% vs 실제 커버리지 33\%의 괴리)을 있는 그대로 드러낸다. 즉 (2) 사용량 데이터는 선행 노출도 지표가 측정하지 못했던 축---"예측된 잠재력 중 실제로 실현된 부분이 얼마인가"---을 직접 관측치로 채워 넣는다. 이것이 정확히 \citet{handa-2025-which-economic-tasks-are}이 스스로를 "이론적 노출지수와 대비되는 실측(observed usage) 접근"이라 규정한 근거이며, \texttt{job\_exposure.csv}가 Eloundou $\beta$를 $WorkUsage_t$로 게이팅하는 방식은 이 관계를 계산식으로 명시한 것이다.

\subsection{(1) 노동시간 데이터 $\leftrightarrow$ AI 적용률(adoption) 개념: 이분법에서 연속적 실현 크기로}

\texttt{06\_AI채택\_실사용\_데이터방법론.tex}가 정리한 채택(adoption) 연구---한국 기업활동조사$\cdot$WPS의 도입 여부 더미, \citet{bick-2026-mind-the-gap-ai-adoption}의 근로자 자기보고 채택률---는 모두 "AI를 도입/사용하는가"라는 \textbf{이분법적} 질문에 응답자가 답하는 구조다. \citet{cheon-2025-ai-exposure-vs-ai-adoption}(전병유$\cdot$신영민, 2025)의 AI Firm Exposure(AIFE)는 이분법을 벗어나긴 하지만, 이 역시 "AI 공급기업의 서비스가 이 직업의 과업을 대상으로 존재하는가"라는 \textbf{공급 측의 존재 여부}를 LLM이 텍스트 매칭으로 판정한 지표이지, 그 서비스가 실제로 그 과업의 소요시간을 얼마나 줄이는지는 측정하지 않는다.

(1) \texttt{human\_only\_time}/\texttt{human\_with\_ai\_time} 데이터는 이 지형에서 한 단계 더 나아간다---"AI를 쓰는가/쓸 수 있는 서비스가 있는가"라는 존재 여부를 넘어, \textbf{그 사용이 실제로 과업 소요시간을 몇 퍼센트 줄이는가}라는 연속적 크기(time savings, 0$\sim$100\%)를 직접 측정한다. 이는 개념적으로 \citet{eloundou-2023-gpts-are-gpts-an-early}의 $\beta$ 등급(0/0.5/1의 3단계 서수)이 "이론상 2배 이상 속도향상이 가능한가"라는 질문에 대한 \textbf{예측된 상한}이었다면, \texttt{human\_only\_time}/\texttt{human\_with\_ai\_time}은 그 예측이 실제 대화에서 \textbf{얼마나 실현되었는지를 사후적으로 재는 연속변수}라는 관계로 요약된다. \citet{tamkin-2025-estimating-ai-productivity-gains-from}가 자신의 연구를 "폭(breadth, 어떤 과업에 AI가 쓰이는가)이 아니라 깊이(depth, 그 과업 내에서 실제로 얼마나 절감되는가)를 포착하는 새 측정"이라 규정한 것이 정확히 이 지점이다.

\subsection{종합: "이론적 상한 $\times$ 실현 정도"라는 공통 구조}

두 데이터를 나란히 놓으면, AEI가 암묵적으로 구현하고 있는 구조는 다음과 같이 요약된다.

\[
\underbrace{\text{선행 노출도 지표(예: Eloundou }\beta\text{)}}_{\text{이론적 상한: 해당 과업이 AI로 처리 가능한가}}
\;\times\;
\underbrace{\text{AEI 사용량 데이터(\S\ref{sec:usage})}}_{\text{실현 여부: 실제로 관측되었는가}}
\;\times\;
\underbrace{\text{AEI 노동시간 데이터(\S\ref{sec:time})}}_{\text{실현 크기: 관측되었다면 얼마나 절감되는가}}
\]

\texttt{job\_exposure.csv}는 이 중 앞의 두 항만을 결합한 것(이론상한$\times$실현여부, 이분법적 게이트)이며, 세 번째 항(실현 크기)까지 결합하는 지표는 AEI 공식 데이터셋에는 아직 존재하지 않는다---\citet{tamkin-2025-estimating-ai-productivity-gains-from}의 시간절감 집계와 \citet{massenkoff-2026-labor-market-impacts-of-ai}의 관측노출은 같은 저장소 내에 있지만 두 논문 모두 서로의 지표를 자신의 공식 안에 아직 직접 대입하지는 않았다. 이는 본 프로젝트가 \texttt{data/proc/merged/}에서 시도해볼 수 있는, AEI 스스로도 아직 완성하지 않은 결합 지점이다.

% ============================================================
\section{결론: \texttt{data/proc/merged/} 설계에 대한 시사점}
\label{sec:conclusion}
% ============================================================

본 해부가 제공하는 가장 구체적인 시사점은 매칭 키의 설계다. AEI 원자료를 한국 노동시장 데이터(한국고용직업분류, 통계청 기업활동조사 등)와 결합하려면, 다음 세 단계의 크로스워크가 필요하다.

\begin{enumerate}[label=(\arabic*)]
  \item \textbf{직업 렌즈 사용 시}: \texttt{onet\_task}/\texttt{soc\_occupation}의 O*NET/SOC 코드 $\to$ (미국 중심 분류체계이므로) 한국고용직업분류(KECO)로의 별도 크로스워크가 필요하며, \texttt{05\_Felten계승연구\_직업분류매칭여부.tex}가 지적했듯 이 크로스워크 자체가 방법론적 선택지가 된다.
  \item \textbf{요청 렌즈 사용 시}: \texttt{request} facet은 직업 분류에 의존하지 않으므로 크로스워크 부담이 적으나, 대신 한국 노동시장 통계(임금$\cdot$고용 등)와 연결할 표준 매핑이 없어 별도 구축이 필요하다.
  \item \textbf{시간 데이터 결합 시}: \texttt{human\_only\_time}/\texttt{human\_with\_ai\_time}을 사용하려면 \S\ref{sec:time-caveat}에서 지적한 단위 검증을 반드시 선행해야 하며, 이를 한국 임금 자료(고용형태별근로실태조사 등)와 결합해 화폐가치화하려면 \citet{fan-2026-aggregate-gains-from-ai}의 LCE 산식을 참고할 수 있다.
\end{enumerate}

이 세 크로스워크 키(O*NET/SOC 코드, 요청군집 코드, 시간단위 검증)를 \texttt{01\_import\_*.do}의 주석에 명시하는 것이, \texttt{CLAUDE.md}가 요구하는 "매칭 키를 log.md 또는 do 파일 주석에 명시" 원칙을 이 데이터에 대해 충족하는 구체적 방법이다.

% ============================================================
\printbibliography[title={참고문헌}]
% ============================================================

\end{document}
